% Part: normal-modal-logic
% Chapter: frame-correspondence
% Section: properties-accessibility

\documentclass[../../../include/open-logic-section]{subfiles}

\begin{document}

\olfileid[id]{nml}{frd}{acc}

\olsection{Sifat Relasi Aksesibilitas}

Banyak !!{formula}s modal ternyata bersifat karakteristik bagi sifat-sifat
sederhana, bahkan sudah dikenal, dari relasi aksesibilitas. Dalam satu arah,
hal ini berarti bahwa setiap model yang mempunyai suatu sifat tertentu membuat
!!{formula} yang bersesuaian (beserta semua instansiasi substitusinya) benar.
Kita mulai dengan lima contoh klasik jenis relasi aksesibilitas dan formula
yang dijamin benar olehnya.

\begin{thm}\ollabel{thm:soundschemas}
  Misalkan $\mModel{M}=\tuple{W,R,V}$ merupakan suatu model. Jika $R$
  mempunyai sifat pada sisi kiri \olref{tab:five}, setiap instansiasi
  !!{formula} pada sisi kanan benar dalam~$\mModel{M}$.
\end{thm}

\begin{table}[t]
    \begin{tabular}{| l || l |}
      \hline
      {\emph{Jika $R$ bersifat \dots}} & {\emph{maka \dots benar dalam~$\mModel{M}$:}} \\
      \hline \hline
      \emph{serial}: $\forall u \exists v Ruv$ & \hbox to.43\textwidth
           {$\Box p \lif \Diamond p$ \hfill (\Ax{D})} \\
      \hline
      \emph{refleksif}: $\forall w Rww$
      & \hbox to.43\textwidth
          {$\Box p \lif p$ \hfill (\Ax{T})} \\
      \hline
      \emph{simetris}: &  \hbox to .43\textwidth
           {$p \lif \Box\Diamond p$ \hfill (\Ax{B})} \\
           $\forall u\forall v(Ruv \lif Rvu)$ & \\
      \hline
      \emph{transitif}: & \hbox to.43\textwidth
           {$\Box p \lif \Box \Box p$ \hfill (\Ax{4})} \\
      $\forall u \forall v \forall w ((Ruv \land Rvw) \lif Ruw)$ & \\
      \hline
      \emph{Euklidean}: & \hbox to .43\textwidth
        {$\Diamond p \lif \Box\Diamond p$ \hfill (\Ax{5})} \\
      $\forall w \forall u \forall v ((Rwu \land Rwv) \lif Ruv)$ &\\
      \hline
    \end{tabular}
    \caption{Lima fakta korespondensi.}
    \ollabel{tab:five}
  \end{table}


\begin{proof}
  Berikut adalah kasus bagi \Ax{B}. Untuk menunjukkan bahwa skema ini benar
  dalam suatu model, kita harus menunjukkan bahwa semua instansiasinya benar
  pada semua dunia dalam model tersebut. Jadi, misalkan
  $!A\lif\Box\Diamond !A$ merupakan suatu instansiasi \Ax{B}, dan misalkan
  $w\in W$ suatu dunia sebarang. Andaikan anteseden~$!A$ benar pada~$w$;
  kita akan menunjukkan bahwa $\Box\Diamond !A$ benar pada~$w$. Jadi, kita
  harus menunjukkan bahwa $\Diamond !A$ benar pada semua~$w'$ yang dapat
  diakses dari~$w$. Bagi sebarang $w'$ sedemikian sehingga $Rww'$, dengan
  menggunakan hipotesis simetri kita juga memperoleh $Rw'w$ (lihat
  \olref{fig:Bsymm}). Karena $\mSat{M}{!A}[w]$, kita mempunyai
  $\mSat{M}{\Diamond !A}[w']$. Karena $w'$ merupakan dunia sebarang sedemikian
  sehingga $Rww'$, kita memperoleh $\mSat{M}{\Box\Diamond !A}[w]$.

  Kasus lainnya ditinggalkan sebagai latihan.
\end{proof}

\begin{prob}
  Lengkapi pembuktian \olref[nml][frd][acc]{thm:soundschemas}.
\end{prob}

\begin{figure}
  \begin{center}
    \begin{tikzpicture}[modal]
      \node[world] (w1) [label={below:$\mSat{{}}{\formula{A}}$\\
          $\mSat{{}}{\Box\Diamond \formula{A}}$}] {$w$};
      \node[world] (w2) [label=below:{$\mSat{{}}{\Diamond \formula{A}}$},
        right=of w1] {$w'$};
      \draw[->, bend left] (w1) to (w2);
      \draw[->, bend left] (w2) to (w1);
    \end{tikzpicture}
  \end{center}
\caption{Argumen berdasarkan simetri.}
\ollabel{fig:Bsymm}
\end{figure}

Perhatikan bahwa implikasi konvers dari \olref{thm:soundschemas} tidak
berlaku: tidak benar bahwa jika suatu model memverifikasi sebuah skema, maka
relasi aksesibilitas model itu mempunyai sifat yang bersesuaian. Dalam kasus
\Ax{T} dan model refleksif, mudah diberikan contoh model tempat \Ax{T} sendiri
gagal: misalkan $W=\{w\}$ dan $V(p)=\emptyset$. Maka $R$ tidak refleksif,
tetapi $\mSat{M}{\Box p}[w]$ dan $\mSat/{M}{p}[w]$. Namun, di sini kita hanya
mempunyai satu instansiasi \Ax{T} yang gagal dalam~$\mModel{M}$;
instansiasi lain, misalnya $\Box\lnot p\lif\lnot p$, benar. Lebih sulit
memberikan contoh tempat \emph{setiap instansiasi substitusi} dari~\Ax{T}
benar dalam~$\mModel{M}$ sementara $\mModel{M}$ tidak refleksif. Namun,
model semacam itu memang ada:

\begin{prop}\ollabel{prop:reflexive}
  Misalkan $\mModel{M}=\tuple{W,R,V}$ merupakan model dengan $W=\{u,v\}$,
  dan dunia $u$ dan~$v$ saling berelasi melalui~$R$; yakni, $Ruv$ dan~$Rvu$.
  Andaikan bagi semua $p$ berlaku $u\in V(p)\Leftrightarrow v\in V(p)$.
  Maka:
  \begin{enumerate}
  \item Bagi semua~$!A$, $\mSat{M}{!A}[u]$ jika dan hanya jika
    $\mSat{M}{!A}[v]$ (gunakan induksi pada~$!A$).
  \item Setiap instansiasi~\Ax{T} benar dalam~$\mModel{M}$.
  \end{enumerate}
  Karena $\mModel{M}$ tidak refleksif (bahkan, model itu
  \emph{irefleksif}), konvers \olref{thm:soundschemas} gagal dalam kasus
  \Ax{T} (argumen serupa dapat diberikan bagi beberapa---meskipun tidak
  semua---skema lain yang disebutkan dalam \olref{thm:soundschemas}).
\end{prop}

\begin{prob}
  Buktikan klaim-klaim dalam \olref[nml][frd][acc]{prop:reflexive}.
\end{prob}

Meskipun kita akan memusatkan perhatian pada lima !!{formula}s klasik \Ax{D},
\Ax{T}, \Ax{B}, \Ax{4}, dan~\Ax{5}, dalam \olref{tab:anotherfive} kita
mencatat beberapa sifat relasi aksesibilitas lainnya. Relasi aksesibilitas~$R$
bersifat fungsional parsial jika dari setiap dunia paling banyak satu dunia
dapat diakses. Jika dari setiap dunia tepat satu dunia dapat diakses, kita
menyebutnya fungsional. (Jadi, relasi fungsional tepat merupakan relasi yang
sekaligus serial dan fungsional parsial.) Relasi ini disebut ``fungsional''
karena relasi aksesibilitas bekerja seperti suatu fungsi (parsial). Suatu
relasi bersifat rapat lemah jika, setiap kali $Ruv$, terdapat suatu~$w$ ``di
antara'' $u$ dan~$v$. Jadi, relasi rapat lemah dalam suatu arti merupakan
kebalikan relasi transitif: dalam relasi transitif, setiap kali $v$ dapat
dicapai dari~$u$ melalui jalan memutar lewat~$w$, $v$ dapat dicapai langsung
dari~$u$; dalam relasi rapat lemah, setiap kali $v$ dapat dicapai langsung
dari~$u$, $v$ juga dapat dicapai melalui jalan memutar lewat suatu~$w$. Suatu
relasi bersifat terarah lemah jika, setiap kali dunia $u$ dan~$v$ dapat
dicapai dari suatu dunia~$w$, terdapat satu dunia~$t$ yang dapat dicapai baik
dari $u$ maupun~$v$---sifat ini kadang-kadang disebut ``sifat wajik'' atau
``konfluensi''.

\begin{table}[t]
    \begin{tabular}{| p{.50\textwidth} || p{.43\textwidth} |}
      \hline
      {\emph{Jika $R$ bersifat \dots}} & {\emph{maka \dots benar dalam~$\mModel{M}$:}} \\
      \hline \hline
      \emph{fungsional parsial}: &
      \multirow{2}{*}{$\Diamond p \lif \Box p$} \\
      $ \forall w \forall u \forall v ((Rwu \land Rwv) \lif u =v )$ & \\
      \hline
      \multirow{2}{*}{\emph{fungsional}: $\forall w \exists v \forall u(Rwu \liff \eq[u][v]) $}
      & \multirow{2}{*}{$\Diamond p \liff \Box p$} \\
      & \\
      \hline
      \emph{rapat lemah}: &  \multirow{2}{*}{$\Box \Box p \lif
        \Box p$} \\
      $\forall u\forall v(Ruv \lif \exists w(Ruw \land Rwv))$ & \\
      \hline
      \emph{terhubung lemah}: &
      \multirow{3}{*}{\hbox to.43\textwidth{$
        \begin{array}{@{}l@{}}
          \Box ((p \land \Box p) \lif q) \lor {} \\
          \qquad \Box ((q \land \Box q) \lif p)
        \end{array}$ \hfill (\Ax{L})}
      } \\
      $\forall w \forall u \forall v ((Rwu \land Rwv) \lif {}$ &\\
      \qquad $(Ruv \lor u=v \lor Rvu))$ & \\
      \hline
      \emph{terarah lemah}: & \multirow{3}{*}{\hbox to .43\textwidth
        {$\Diamond\Box p \lif \Box\Diamond p$\hfill (\Ax{G})}} \\
      $\forall w \forall u \forall v ((Rwu \land Rwv) \lif {}$ & \\
      \qquad $\exists t (Rut \land Rvt))$ & \\
      \hline
    \end{tabular}
    \caption{Lima fakta korespondensi lainnya.}
    \ollabel{tab:anotherfive}
  \end{table}

\begin{prob}
  Misalkan $\mModel{M}=\tuple{W,R,V}$ merupakan suatu model. Tunjukkan bahwa
  jika $R$ memenuhi sifat-sifat pada sisi kiri
  \olref[nml][frd][acc]{tab:anotherfive}, setiap instansiasi formula yang
  bersesuaian pada sisi kanan benar dalam~$\mModel{M}$.
\end{prob}

\end{document}
